\begin{center}
	\Large{\textbf{EXECUTIVE SUMMARY}}
\end{center}
\vspace{1em}
\normalsize
\justifying
\hspace{1cm}The QNA-Auth prototype is a proof of concept designed to allow devices that use noise-based verification through their camera or microphone. The system's construction has a formal, conservative boundary. We do not consider learned embeddings to be digital certificates or proof of cryptographic identity; rather, we prefer to think of them as biometrically-like feature template images that can support continuous similarity comparisons under restricted conditions. This delineation of the engineering work is important for continuing to provide an unbiased engineering basis while permitting an informative study of device verification.


\hspace{1cm}The overall implementation consists of signal processing techniques, machine learning algorithms, and a security oriented system design. Raw samples from the cameras and microphones are converted into a canonical representation of the feature that includes statistics, frequencies, entropies, autocorrelations, and complexity. A Siamese-style embedding network maps the feature vectors into a common similarity measure, such that the same device's and combined representations can be compared throughout both registration and authentication. While making respect to single threshold values, run-time decisions are based on multiple thresholds through the use of weighted multi-source fusion, confidence bands, source-profile guards, and impostor-margin checks, thus providing a means to reduce decision brittle-ness.

\hspace{1cm}The QNA-Auth system has more than just the core of QNA but also a FastAPI backend, using SQLite for persistence, log audits, React for the frontend, diagnostic scripts; a nonce-based challenge / response layer that is made more secure through HKDF and a server-side secret. It further includes evaluation artifacts, reviewer notes, and demo documentation for the entire system. All components provide a thesis project that shows the full end-to-end, technically and functionally sound, device verification pipeline, while being quite explicit about the limitations of the dataset, environmental drift, and the need for additional validation in the future.
